\section{Beskriv problemet}

\subsection{Omfang}

Bergen Traffic Predictor estimerer hvor mange kjøretøy som passerer ett bestemt
trafikkregistreringspunkt i løpet av neste time. Punktet er Danmarks plass med ID
\texttt{14630V805616}. Hver modellrad representerer et tidspunkt $t$, mens
\emph{target} er observert trafikkvolum ved $t+1$. Oppgaven er dermed
\emph{supervised regression} og en kortsiktig tidsserieprognose.

Prosjektets opprinnelige hypotese var at historiske værdata fra MET Norways Frost
API kunne forbedre prediksjonene. Værvariantene tapte i kronologisk validering, så
sluttmodellen ble værfri som følge av eksperimentene.

En videreutviklet løsning kan gi trafikkplanleggere og analytikere et anslag for
belastningen ved punktet. Den manuelle referansen er trafikken på samme klokkeslett
én uke tidligere, eventuelt justert skjønnsmessig for helg eller ferie. Dette er
prosjektets baseline. Maskinlæring er relevant fordi trafikken varierer med time,
ukedag og kalenderhendelser, og fordi trafikkverdier fra flere tidligere
tidspunkter kan samvirke ikke-lineært. Modellen mangler likevel informasjon om blant annet ulykker,
vegarbeid og arrangementer.

Prognosen kan brukes til å oppdage avvik fra et normalt ukemønster eller som ett av
flere grunnlag for planlegging. Studentprosjektet har imidlertid ikke undersøkt
faktisk beslutningsbruk, brukertall eller økonomisk gevinst. Løsningen demonstrerer
derfor teknisk gjennomførbarhet, ikke dokumentert nytte eller lønnsomhet i drift.

Avgrensningen til ett punkt og én times horisont gjør dataløpet og evalueringen
entydig. Modellen sier ikke noe om reisetid, kølengde eller retning og kan ikke
overføres direkte til andre målepunkter med andre trafikkmønstre.

Figur~\ref{fig:pipeline} viser systemgrensen. Traffic Data API leverer aggregerte
timesmålinger med kvalitetsparametere \cite{vegvesen-traffic-api}. Prosjektet
normaliserer serien, bygger modellvariabler fra eksakte tidspunkter, trener modellen
og bruker den i Streamlit. Trening og app deler feature-kode, mens datakvalitet og
publiseringsforsinkelse avgjør om en prognose kan lages.

\begin{figure}[htbp]
  \centering
  \resizebox{\textwidth}{!}{\begin{tikzpicture}[
  node distance=5.5mm and 5.5mm,
  stage/.style={draw=black!55, rounded corners=1.5pt, fill=black!3,
    minimum height=10mm, text width=27mm, align=center, font=\small},
  arrow/.style={-{Latex[length=2mm]}, semithick, draw=black!70}
]
\node[stage] (api) {Traffic Data API\\timevolum};
\node[stage, right=of api] (prepare) {Validering og\\komplett timeakse};
\node[stage, right=of prepare] (features) {Kalender og eksakte\\lag på tidspunkt};
\node[stage, right=of features] (model) {Frosset HGB\\og metadata};
\node[stage, right=of model] (app) {Lokal\\Streamlit-app};
\draw[arrow] (api) -- (prepare);
\draw[arrow] (prepare) -- (features);
\draw[arrow] (features) -- (model);
\draw[arrow] (model) -- (app);
\end{tikzpicture}
}
  \caption{Implementert dataløp fra offentlige trafikkdata til lokal demonstrasjon.}
  \label{fig:pipeline}
\end{figure}
\FloatBarrier

Prosjektet regnes som vellykket dersom en lekkasjesikker modell slår
forrige-uke-baselinen på en senere, urørt periode, kan bygges reproduserbart og gir
enten prognose eller tydelig feilmelding. Flerårige timedata og HGB-modellen kan
håndteres på en vanlig bærbar datamaskin. Den viktigste ressursen er
utviklingsarbeid til datavalidering, lekkasjekontroll, evaluering og integrasjon.
En produksjonsløsning ville i tillegg krevd driftsansvar, overvåking og rutiner for
godkjenning av nye modellversjoner.

\subsection{Metrikker}

Hovedmetrikken er gjennomsnittlig absolutt feil, MAE,

\begin{equation}
  \operatorname{MAE}=\frac{1}{n}\sum_{i=1}^{n}|y_i-\hat{y}_i|.
\end{equation}

MAE måles i kjøretøy per time og kan tolkes direkte. RMSE rapporteres som støttemål
fordi kvadreringen gir større utslag for enkelte store feil:

\begin{equation}
  \operatorname{RMSE}=\sqrt{\frac{1}{n}\sum_{i=1}^{n}(y_i-\hat{y}_i)^2}.
\end{equation}

Modell og baseline måles på de samme prognoseradene. Vi rapporterer også antall
prediksjoner og feil etter time og ukedag. Felles rader hindrer at ulik
tilgjengelighet av historiske trafikkverdier gir forskjellige evalueringsutvalg.
Programvaren testes for reproduserbar bygging, modellasting, manglende historiske
trafikkverdier, API-feil og tidssoner. Svartid,
oppetid og brukertilfredshet er ikke målt i den lokale demonstrasjonen.
